\chapter{Discussion}
\label{ch:discussion}

The empirical evaluation presented in Chapter~\ref{ch:evaluation_results}
demonstrated the feasibility, structural concordance, and computational
scalability of OmniDeps across 156 real-world software repositories and five
programming languages. However, interpreting these quantitative findings
requires stepping back from raw similarity metrics and runtime curves to
examine the broader architectural, methodological, and practical implications
of a language-agnostic approach to static analysis.

This chapter synthesizes the empirical outcomes into a cohesive discussion
structured around four primary themes:
\begin{enumerate}
  \item \textbf{Key Empirical Findings and Paradigm Dynamics}: What do the
    comparative metrics reveal about dependency extraction across divergent
    type systems and compilation models? Why do certain language idioms
    exhibit high structural consensus while others diverge?
  \item \textbf{Architectural Efficacy Assessment}: Was the design decision
    to replace compiler-coupled frontends with lightweight, AST-querying
    parsers worthwhile? What are the exact trade-offs between build-system
    immunity and deep compiler-level type resolution?
  \item \textbf{Practical Implications for Software Engineering Tooling}:
    How can software engineering practitioners, tool authors, and CI/CD
    architects leverage these insights for continuous architectural
    governance and smell detection?
  \item \textbf{Threats to Validity}: What methodological, construct, and
    empirical limitations govern the interpretation and generalizability of
    our results?
\end{enumerate}

% ==============================================================================
% 1. KEY EMPIRICAL FINDINGS
% ==============================================================================
\section{Key Empirical Findings and Paradigm Dynamics}
\label{sec:discussion_findings}

Our empirical findings uncover fundamental differences in how programming
language paradigms interact with syntax-directed static dependency extraction.
Rather than reflecting uniform performance across the corpus, the results
highlight distinct extraction dynamics across object-oriented, dynamic,
procedural, and modern systems languages.

\subsection{High Semantic Parity in Object-Oriented Ecosystems}
\label{subsec:discussion_oo_parity}

The most prominent empirical finding in object-oriented environments (Java) is
the near-complete convergence of architectural entities and the high degree of
relational concordance between OmniDeps and compiler-backed analyzers. As
reported in Section~\ref{subsec:rq2_cross_tool}, OmniDeps achieved an entity
inferiority metric of just $\text{Inf}_{\text{nodes}} = 0.021$ against
Depends and $0.003$ against Arcan, while edge Jaccard similarity converged to
$\approx 0.58$--$0.60$ across all three baseline tools.

This convergence validates a foundational architectural hypothesis: \textbf{in
explicitly typed, modular object-oriented languages, Concrete Syntax Trees
contain virtually all structural information required for architectural
dependency modeling}. Java source files exhibit an unambiguous syntactic
isomorphism:
\begin{itemize}[noitemsep]
  \item Source files map cleanly to canonical package hierarchies via explicit
    \texttt{package} declarations.
  \item Types (classes, interfaces, enums, records) are explicitly declared
    with declared subtyping relations (\texttt{extends}, \texttt{implements}).
  \item Dependency targets are explicitly imported via \texttt{import}
    statements or qualified inline references.
\end{itemize}

Because compilation units in Java cannot inject arbitrary global preprocessor
macros or alias global namespaces dynamically, OmniDeps' two-phase algebraic
query model---resolving lexical scopes upward via \texttt{Find} and descending
inheritance graphs via \texttt{Extract}---successfully resolves approximately
$98\%$ of compiler-resolved architectural entities without requiring access to
external JAR dependencies or a live Java Virtual Machine (JVM).

Furthermore, the detailed categorical discordance analysis
(Section~\ref{subsec:rq3_discordance}) confirmed that containment hierarchies
exhibit over $98.2\%$ pairwise agreement. The remaining relational divergence
is primarily driven by symmetric call coupling and member-level field
references rather than missed architectural components. This confirms that
for high-level architectural static analysis (e.g., architectural smell
detection, module cohesion, package cycle identification), syntax-directed
querying matches the fidelity of complete compiler symbol tables.

\subsection{The Resolution Ceiling in Dynamically Typed Languages}
\label{subsec:discussion_dynamic_languages}

In contrast to object-oriented environments, evaluating Python codebases
uncovered a pronounced structural dichotomy: while entity identification
remains robust ($J_{\text{nodes}} = 0.740$), call-level dependency alignment
exhibits substantial divergence ($J_{\text{edges}} = 0.296$).

This discrepancy is not an implementation artifact, but reflects a
fundamental \textbf{information-theoretic resolution ceiling inherent to
static analysis of dynamically typed languages}. In Python:
\begin{itemize}
  \item \textbf{Absence of Static Type Annotations}: In idiomatic Python,
    function arguments and variables frequently omit type signatures (e.g.,
    \texttt{def process(handler): handler.dispatch()}). Without dynamic
    execution traces or whole-program flow-sensitive type inference, a
    syntax-directed static analyzer cannot identify which class instantiates
    \texttt{handler}.
  \item \textbf{Dynamic Module Mutation}: Python allows runtime modification
    of module symbol tables via \texttt{sys.modules}, dynamic imports
    (\texttt{importlib.import\_module}), wildcard imports (\texttt{from foo
    import *}) evaluated at runtime, and dynamic re-exports controlled by
    \texttt{\_\_all\_\_}.
  \item \textbf{Late Binding and Duck Typing}: Method invocations are
    dispatched dynamically by attribute name lookup (\texttt{\_\_getattr\_\_},
    \texttt{\_\_getattribute\_\_}) rather than static virtual method tables.
\end{itemize}

While Depends attempts to mitigate this by maintaining speculative symbol
bindings across visited files, both OmniDeps and Depends inevitably arrive at
divergent heuristics when resolving ambiguous method calls. This empirical
result demonstrates that without incorporating gradual type annotations (PEP
484 \cite{pep484}) or constraint-based abstract interpretation, purely
syntactic static analysis in dynamic languages cannot achieve the edge
recall levels observed in statically typed ecosystems.

\subsection{The Procedural Dilemma in C and C++: Headers, Macros, and Translation Units}
\label{subsec:discussion_procedural_dilemma}

The empirical evaluation of C and C++ repositories revealed the most intricate
architectural trade-offs of the entire study, governed by two intertwined
factors: preprocessor macro expansion and header inlining.

\paragraph{Preprocessor Macro Masking in Raw Sources.}
In raw, unprocessed C and C++ source code, extensive use of preprocessor
macros (\texttt{\#define}) fundamentally obscures the underlying syntax.
Macros are routinely used to declare function wrappers, synthesize struct
typedefs, declare interface dispatch tables, and implement conditional
compilation (\texttt{\#ifdef}). As shown in Table~\ref{tab:results_summary_rq2},
analyzing raw C code yields a modest edge Jaccard similarity of $0.115$. Tree-sitter
parsers, encountering macro invocations that do not conform to standard C
grammar, frequently classify complex macro statements as generic call
expressions or syntax errors, causing referenced structs and typedefs to remain
unlinked.

\paragraph{Translation-Unit Expansion vs. Header Duplication.}
When translation units are preprocessed via the compiler frontend
(\texttt{clang -E}), macros are fully expanded, conditional directives are
evaluated, and all included header files are inlined directly into the source
stream. As documented in Section~\ref{subsec:rq2_cross_tool}, this produces a
substantial boost in structural recall: in C++, edge Jaccard agreement doubles
from $0.081$ to $0.165$ (reaching $0.459$ on \texttt{re2} and $0.320$ on
\texttt{ninja}), while node inferiority drops significantly.

However, translation-unit preprocessing introduces an acute architectural
distortion: \textbf{header file inlining duplicates shared type declarations
across every single compilation unit}. A common struct declared in a core header
(e.g., \texttt{redis.h} or \texttt{sqliteInt.h}) is replicated dozens or
hundreds of times across the preprocessed translation units. This artifact
artificially inflates node superiority ($\text{Sup}_{\text{nodes}} = 7.22$ in C
and $2.14$ in C++), because OmniDeps extracts each inlined instance as an
autonomous entity within its local translation unit.

\paragraph{The Incomplete Type Trap in Header Stripping.}
To address this duplication, our exploratory investigation evaluated
stripping inlined header declarations using compiler linemarkers. However,
this experiment exposed a fatal flaw in compiler-dependent baseline tools:
removing header definitions leaves source files with incomplete type
declarations (e.g., \texttt{struct Point;}). While OmniDeps' query engine
tolerates incomplete forward declarations, compiler frontends (such as Eclipse
CDT within Depends) require complete struct layouts to resolve member offsets
and pointer dereferences. Consequently, CDT crashed with syntax and unresolved
symbol exceptions across nearly every stripped project, proving that
compiler-based frontends cannot function without complete header inlining.

These empirical observations lead to a decisive conclusion: \textbf{in
procedural C and C++, translation units cannot be treated as autonomous
modular components}. Bridging the gap between physical files and logical
architecture requires an explicit, post-parse link-time deduplication layer
that unifies identical struct declarations across translation units while
preserving call edges.

\subsection{Modern Modularity and Trait Systems in Rust}
\label{subsec:discussion_rust_modularity}

The evaluation of Rust repositories provides an instructive counterpoint to C
and C++. Rust combines low-level systems programming performance with a modern,
strictly enforced modularity model:
\begin{itemize}[noitemsep]
  \item Module hierarchies map strictly to the filesystem hierarchy
    (\texttt{mod.rs} or matching file basenames).
  \item Dependencies across crates and submodules are explicitly declared via
    canonical \texttt{use} paths.
  \item Behavioral polymorphism is governed by explicit trait implementation
    blocks (\texttt{impl Trait for Struct}) rather than header files or
    preprocessor macros.
\end{itemize}

OmniDeps successfully analyzed 100\% of the 32 Rust repositories, sustaining a
throughput of over 97,000 LOC/s. Because Rust's grammar provides unambiguous
syntactic demarcations for struct definitions, enum variants, and trait
implementations, the algebraic query engine resolves cross-module type
references with exceptional determinism. This confirms that modern systems
languages engineered with clean module systems allow CST-based static analysis
to achieve compiler-level fidelity while completely bypassing compiler
toolchain overhead.

% ==============================================================================
% 2. ARCHITECTURAL EFFICACY ASSESSMENT
% ==============================================================================
\section{Architectural Efficacy: Was the Agnostic Approach Worth It?}
\label{sec:discussion_efficacy}

A foundational question guiding this dissertation is whether developing a
language-agnostic, Concrete Syntax Tree-querying static analyzer is
architecturally and practically worthwhile compared to maintaining bespoke,
compiler-coupled frontends.

Based on our empirical findings across 156 repositories, the answer is a
resounding \textbf{affirmative}. While compiler-coupled frontends offer marginal
precision advantages in specialized scenarios (such as complex C++ template
metaprogramming), the benefits of the decoupled AST-query architecture
overwhelmingly outweigh these edge cases across three critical dimensions:
operational feasibility, computational throughput, and engineering economy.

\subsection{The Build-System Decoupling Dividend}
\label{subsec:discussion_build_decoupling}

The single most consequential finding of our comparative evaluation is the
dramatic disparity in operational completion rates between OmniDeps and
traditional static analyzers:
\begin{itemize}
  \item \textbf{OmniDeps}: \textbf{99.4\% completion} (155/156 projects analyzed
    successfully; 100\% in C, Java, Python, and Rust).
  \item \textbf{Depends}: \textbf{92.7\% completion} (115/124 projects).
  \item \textbf{Arcan}: \textbf{54.0\% completion} (34/63 projects).
  \item \textbf{Entangle}: \textbf{29.7\% completion} (19/64 projects).
\end{itemize}

In empirical software engineering and industrial static analysis, the
requirement that an analyzer must fully integrate with a project's build
system is a catastrophic point of failure. Compiler-coupled tools (such as
Arcan, which relies on compiler frontends, or Clang LibTooling-based systems)
require:
\begin{enumerate}[noitemsep]
  \item An exact, matching compiler toolchain and SDK installed on the host.
  \item Fully resolved external dependencies (e.g., all Maven/Gradle JARs
    downloaded, C/C++ third-party libraries installed in system search paths).
  \item A valid, synchronized build configuration file (e.g.,
    \texttt{compile\_commands.json}, \texttt{pom.xml}, \texttt{CMakeLists.txt}).
\end{enumerate}

When evaluating real-world open-source repositories, these prerequisites fail
regularly: build scripts rely on private repository credentials, deprecated
build plugins fail on modern environments, platform-specific header paths are
missing, or submodules are uninitialized. As a result, Arcan failed to analyze
nearly half of the target corpus (46.0\% failure rate).

By contrast, OmniDeps operates under the foundational premise that
\textbf{the source code compiles successfully within its native environment},
deliberately decoupling syntax analysis from build-system bootstrapping.
OmniDeps requires nothing more than raw source trees on disk. For empirical
researchers analyzing hundreds of historical releases, or software architects
evaluating multi-language enterprise monorepos, \textbf{an analyzer that
successfully extracts 98\% of dependencies with zero configuration is
infinitely more valuable than an analyzer that achieves 100\% theoretical
precision but fails to run on 46\% of codebases}.

\subsection{Precision vs. Scalability: The Pragmatic Trade-off}
\label{subsec:discussion_precision_tradeoff}

Traditional compiler frontends dedicate immense computational resources to
whole-program semantic analysis: constructing control-flow graphs, solving
interprocedural points-to constraints, and validating type-system correctness.
However, for the purposes of \textbf{software architecture recovery and
dependency modeling}, this level of verification is largely redundant. The
source code has already been verified and compiled by the developer's native
toolchain.

OmniDeps exploits this insight by replacing multi-pass compiler type checking
with a single-pass CST extraction followed by an algebraic Scope Tree query.
The performance dividends of this architectural trade-off are profound:
\begin{itemize}
  \item \textbf{Throughput}: OmniDeps sustains a mean throughput of
    \textbf{74,445 lines of code per second}, analyzing large codebases like
    \texttt{redis} or \texttt{git} in under three seconds.
  \item \textbf{Speedup}: OmniDeps achieves average speedups of \textbf{13.2x
    over Depends}, \textbf{48.2x over Arcan}, and \textbf{4.0x over Entangle}.
\end{itemize}

What is sacrificed in exchange for this $48\times$ speedup? As our Java
experiments demonstrated, OmniDeps incurs an entity omission rate of just
$2.1\%$ against Depends and $0.3\%$ against Arcan. For macro-architectural
analyses---such as detecting Cyclic Dependencies, Hub-Like Components, or
Unstable Interfaces---a 2\% difference in fine-grained call edges does not
alter the identified architectural topology. Consequently, the trade-off
between exhaustive compiler verification and high-throughput AST querying is
pragmatically optimal for architectural analysis.

\subsection{Engineering Economy and Polyglot Extensibility}
\label{subsec:discussion_engineering_economy}

Maintaining a multi-language static analysis framework built on bespoke
compiler frontends imposes an unsustainable engineering burden. Supporting five
distinct languages via compiler plugins requires maintaining active bindings
to:
\begin{itemize}[noitemsep]
  \item Clang LibTooling for C and C++ (C++ API, strict LLVM version pinning).
  \item Eclipse JDT or javac compiler plugins for Java (JVM runtime).
  \item Python's native \texttt{ast} or \texttt{libcst} module (Python runtime).
  \item \texttt{rustc\_driver} or \texttt{ra\_ap\_syntax} for Rust (unstable
    internal compiler APIs).
\end{itemize}

Each of these compiler APIs evolves independently, introduces incompatible AST
abstractions, and requires dedicated runtime environments.

In contrast, OmniDeps establishes an elegant, unified engineering footprint.
By utilizing Tree-sitter, language frontends are defined entirely through
\textbf{declarative S-expression syntax queries} mapping concrete CST patterns
to the strongly typed Universal Intermediate Representation. Adding support
for a new programming language (e.g., Go, C\#, or TypeScript) requires writing
a set of query definitions; the underlying algebraic query engine, Scope Tree
indexing, and graph constructor remain completely unchanged and language-agnostic.
This engineering economy ensures that OmniDeps can be maintained and extended
by small research or engineering teams without polyglot compiler expertise.

% ==============================================================================
% 3. PRACTICAL IMPLICATIONS
% ==============================================================================
\section{Practical Implications for Software Engineering and Tool Design}
\label{sec:discussion_implications}

The findings of this thesis have immediate practical ramifications for software
engineering practitioners, architectural tool vendors, and development
workflows:

\subsection{Continuous Integration and Pre-Commit Architectural Auditing}
\label{subsec:implications_ci_cd}

Traditional architectural smell detection and dependency auditing tools are
frequently relegated to nightly batch jobs or sprint-level audits because their
execution times range from tens of minutes to hours on enterprise repositories.
Because OmniDeps analyzes systems comprising hundreds of thousands of lines of
code in mere seconds, \textbf{architectural static analysis can now be shifted
left into Continuous Integration (CI/CD) pipelines and pre-commit hooks}.

Development teams can enforce architectural governance rules (such as layer
violations, illegal cross-module dependencies, or cyclic component coupling) as
mandatory pull-request checks, catching structural decay before it is merged
into production branches.

\subsection{Universal Ingestion Engine for Architectural Smell Detectors}
\label{subsec:implications_smell_detection}

State-of-the-art architectural analysis platforms such as \textit{Arcan}
\cite{fontana_arcan_2017}, \textit{DV8} \cite{xiao_design_2014}, and
\textit{SonarQube} expend substantial development effort engineering and
maintaining language-specific code parsers.

Our results demonstrate that OmniDeps' Universal IR and normalized graph
harmonization format can serve as a \textbf{universal ingestion frontend} for
these downstream analysis engines. By interfacing with OmniDeps via standardized
graph exports (such as GraphML or JSON), smell detection tools can instantly
expand their language support to C, C++, Java, Python, and Rust while shedding
their dependency on fragile compiler toolchains.

\subsection{Polyglot Monorepo Governance}
\label{subsec:implications_monorepos}

Modern enterprise software systems increasingly adopt polyglot monorepo
architectures, where microservices, frontend applications, and shared libraries
written in divergent languages coexist within a single version-controlled
repository. Traditional static analyzers operate strictly within single-language
silos, leaving cross-language interactions unmodeled.

Because OmniDeps represents software components in a unified, language-neutral
Intermediate Representation with canonical qualified naming, it lays the
technical foundation for holistic polyglot dependency extraction, enabling
practitioners to trace architectural relationships across language boundaries
within a single unified multigraph.

% ==============================================================================
% 4. THREATS TO VALIDITY
% ==============================================================================
\section{Threats to Validity}
\label{sec:threats_to_validity}

To ensure sound interpretation and generalizability of the experimental
outcomes, potential threats to validity must be recognized and systematically
mitigated across construct, internal, and external validity dimensions.

\subsection{Construct Validity}
\label{subsec:construct_validity}

Threats to construct validity concern whether the operationalized evaluation
metrics and graph transformations genuinely reflect architectural dependency
extraction fidelity rather than superficial artifacts of representation.

\begin{itemize}[leftmargin=2.5em]
  \item \textbf{Metamodel Harmonization and Semantic Lifting}:
    Divergent static analysis tools represent software structure at
    fundamentally different granularities. A fine-grained tool such as
    \textit{Depends} captures AST-level function calls and local field
    references, whereas architectural smell detectors such as \textit{Arcan}
    extract class-to-class and package-level couplings. We mitigated this
    disparity through the formal Semantic Lifting Operator ($\lambda$,
    Section~\ref{subsec:semantic_lifting}), which lifts member-level
    interactions to enclosing architectural types ($V_{\text{arch}}$) and
    suppresses intra-component self-loops. Furthermore, we defined Intermediate
    Package Synthesis and Explicit Containment Closure
    (Section~\ref{subsec:structural_closure}) to eliminate artificial
    penalization arising from compound node nesting versus explicit containment
    edges.

  \item \textbf{Ontological Granularity Gaps in Procedural Codebases}:
    In procedural languages (C and C++), analyzers may model architecture at the
    level of physical translation units (file-to-file graphs) rather than
    user-defined data structures (type-to-type graphs). As analyzed in
    Section~\ref{subsubsec:arcan_c_cpp_divergence}, comparing Arcan's file-level
    graphs with OmniDeps' type-centric graphs yields disjoint vertex spaces
    ($V_{\text{Arcan}} \cap V_{\text{OmniDeps}} = \emptyset$). Rather than
    reporting artificially deflated zero scores, we identified this as an
    ontological mismatch and explicitly marked Arcan as non-comparable for C and
    C++, preserving metric construct integrity.

  \item \textbf{Metric Sensitivity and Smoothing}:
    Graph similarity is evaluated using the Laplace-smoothed Jaccard index
    ($J_{\text{nodes}}$ and $J_{\text{edges}}$,
    Section~\ref{subsec:jaccard_similarity}). For sparse submodules or
    repositories where no subtype edges exist, Laplace smoothing prevents
    indeterminate forms ($\frac{0}{0}$) without biasing results on large-scale
    systems where $|V| \gg 1$ and $|E| \gg 1$. In addition, directional
    superiority ($\text{Sup}$) and inferiority ($\text{Inf}$) indices
    (Section~\ref{subsec:discrepancy_indices}) ensure that asymmetric
    over-approximation and omission are quantified independently.
\end{itemize}

\subsection{Internal Validity}
\label{subsec:internal_validity}

Threats to internal validity concern potential confounding factors in the
experimental implementation, preprocessing pipeline, or measurement execution.

\begin{itemize}[leftmargin=2.5em]
  \item \textbf{Translation-Unit Preprocessing and Header Inlining}:
    In procedural and object-oriented C/C++, macro expansion and header file
    inclusion represent a critical trade-off. In raw source files, unexpanded
    conditional macros and missing \texttt{\#include} paths cause unlinked
    types and unresolved function invocations. To evaluate the impact of
    compiler preprocessing, we implemented an automated translation-unit (TU)
    preprocessor invoking \texttt{clang -E} (filtering system headers from
    \texttt{/usr/include}). While TU preprocessing eliminates unlinked macros
    and increases node recall (reducing $\text{Inf}_{\text{nodes}}$ to $\approx
    0.05$--$0.15$ on key projects), inlining shared header files duplicates
    type definitions across every translation unit that includes them, inflating
    apparent node superiority ($\text{Sup}_{\text{nodes}}$). Conversely,
    post-preprocessing stripping of inlined headers via compiler linemarkers
    breaks compiler-based frontends (such as Eclipse CDT in Depends), which
    require complete type definitions within each compilation unit to resolve
    symbols. Consequently, we report both raw and preprocessed results to
    isolate the impact of header inlining.

  \item \textbf{Dynamic Features and Reflection}:
    In dynamically typed languages like Python, static analysis inherently cannot
    observe runtime reflection, dynamic dispatch via \texttt{getattr}, or runtime
    monkey-patching. This creates an expected ceiling on edge recall in dynamic
    languages, as static AST queries cannot infer types that exist solely during
    interpreter runtime.

  \item \textbf{Benchmarking and Execution Measurement Rigor}:
    Comparing native compiled binaries (written in Rust) against JVM-based
    analyzers (Arcan and Depends) introduces potential execution bias due to JVM
    bytecode warmup, Just-In-Time (JIT) compilation, and garbage collection
    pauses. To mitigate runtime variance, each benchmark iteration was executed
    within isolated POSIX process sessions (\texttt{setsid}) with clean subprocess
    tree termination via $\texttt{SIGKILL}$, CPU affinity pinning, and a hard
    timeout ($T_{\text{max}} = 30$ minutes). Timings were captured across repeated
    sampling iterations, separating process startup from steady-state execution.
\end{itemize}

\subsection{External Validity}
\label{subsec:external_validity}

Threats to external validity concern the generalizability of our findings
across diverse software ecosystems, paradigms, and programming languages.

\begin{itemize}[leftmargin=2.5em]
  \item \textbf{Corpus Scale and Heterogeneity}:
    To avoid conclusions biased toward specific project structures, our
    empirical evaluation dataset encompasses 156 open-source software projects
    spanning five distinct languages (31 C, 31 C++, 32 Java, 30 Python, and 32
    Rust repositories) and aggregating over 45 million lines of code.

  \item \textbf{Architectural and Domain Diversity}:
    The evaluated corpus includes low-level systems and operating systems
    tooling (\texttt{git}, \texttt{tmux}, \texttt{libuv}), high-throughput
    databases and storage engines (\texttt{redis}, \texttt{sqlite},
    \texttt{ClickHouse}, \texttt{duckdb}), distributed enterprise platforms
    (\texttt{elasticsearch}, \texttt{flink}), game engines (\texttt{godot}),
    networking microframeworks (\texttt{fastapi}, \texttt{flask},
    \texttt{actix-web}), and developer toolchains (\texttt{ripgrep},
    \texttt{cargo}, \texttt{ninja}). This breadth ensures that evaluation results
    reflect real-world engineering diversity rather than narrow domain-specific
    idioms.
\end{itemize}

\section{Chapter Summary}
\label{sec:discussion_summary}

This chapter synthesized the empirical results into broader architectural and
methodological insights. We analyzed the paradigm dynamics that govern
dependency extraction, explaining the high semantic parity in Java, the
inherent resolution ceiling in dynamic Python, the translation-unit trade-offs
in C/C++, and the clean modularity of Rust. We critically evaluated the
architectural efficacy of the language-agnostic approach, demonstrating that
decoupling from the native build toolchain provides overwhelming advantages in
operational feasibility (99.4\% vs. 54.0\%) and computational throughput
(74,445 LOC/s), far outweighing minor precision trade-offs. Finally, we
highlighted the practical implications for continuous integration and
polyglot repository governance, and formally framed the threats to validity that
bound these conclusions.
